The current evaluation framing is intentionally conservative. It treats the simulator as the primary source of truth and avoids overstating what the present evidence can support. This framing should not be misinterpreted as uncertainty about all findings. What the evidence establishes is concrete: within the proposed architecture, a coupled memory--weighting--reflection loop can be serialized, rerun, and evaluated in a traceable way, and that loop produces measurable but uneven gains inside the local simulator.

First, the strongest aspect of the evidence is the system's internal consistency. The same structured plan objects feed the generator, simulator, result serializer, and DoDAF/C2SIM exporters. This makes it possible to trace reported numbers back to stored JSON files and to inspect the logic that produced them. For a rigorous application-oriented study, such traceability is more valuable than broader but weakly grounded claims.

Second, the selected evidence suggests that the interaction between memory augmentation, adaptive weighting, and reflection is useful within the current simulator, but also that the gains remain scenario-dependent. The full pipeline is best in the chosen single-scenario ablation, leading four of the five headline metrics and four of the five stages, and exceeds the pure baseline in four of five scenarios of the refreshed cross-scenario suite (RTX 4090 D). At the same time, the stage-level evidence shows that the gains are not perfectly uniform. As emphasized by the single-scenario stage table and stage-profile figure, \texttt{breach} remains the most persistent bottleneck (60.88 under Full), and the \texttt{sustain} stage is led by the Hope-only variant (70.09 vs.\ 66.57). The latter observation is consistent with the weight-competition analysis: \texttt{breach} and \texttt{sustain} both draw on the \texttt{protection} capability dimension, and the scalarized reward cannot fully resolve their competition. The BottleneckDetector's dynamic boost directs a larger share of the adaptation budget toward the currently weakest stage, producing the largest Full-vs-next-best gaps in \texttt{disrupt} (+4.4 points) and \texttt{breach} (+2.3 points).

The newly introduced \texttt{BottleneckDetector} module addresses this diagnostic gap by making bottleneck identification explicit and dynamic. The detector's output confirms that breach is the active bottleneck in every iteration of the single-scenario run, with stable severity scores around 0.053--0.054. The dynamic boost vector it produces directs additional emphasis to mobility, fires, and protection---the three capabilities that feed the breach stage---in proportion to the measured severity. While the dynamic boost alone does not eliminate the breach bottleneck (the stage-score evidence shows breach still lagging), it provides a principled mechanism for the controller to concentrate its adjustment budget on the most urgent stage, replacing the original static priority ordering that always privileged disrupt and breach equally. The weight-competition analysis reveals a structural reason for breach's persistence: breach and sustain share the \texttt{protection} capability dimension with equal weight (0.25), creating a tension that a scalarized reward cannot fully resolve. This finding connects the pipeline's behavior to the broader multi-task optimization literature, where gradient conflicts between competing objectives are a known challenge \citep{yu2020pcgrad}.

The memory-conditioned Hope coupling (Section~\ref{sec:memory-coupling}) is a further step toward tighter component interaction. By letting retrieved case metrics bias the fast-weight initialisation, the pipeline no longer treats memory and adaptive weighting as purely serial modules. The correction is intentionally lightweight (clamped to $[-0.18, 0.28]$, blended at 0.30 weight) so that it nudges rather than overrides the adapter's scenario-conditioned output. The practical effect is that when the Memento module retrieves a case that failed badly at breach, the controller enters the next planning round with a slight upward bias on the breach-relevant capability dimensions. This coupling does not require retraining the adapter or modifying the case bank schema, and it means the pipeline's two auxiliary modules now share information through a common capability representation.

A noteworthy observation from the cross-scenario Reflection-only experiments reinforces the case for tight coupling. Under seed 7, Reflection-only underperforms Pure LLM in all five generalization scenes (mean deficit $-2.72$ points). However, a multi-seed replication across three seeds $\{7, 42, 123\}$ on all five scenes reveals that this negative effect is seed-dependent: the mean Reflection-only MS is 63.23 (seed 7, clearly worse than Pure), 65.80 (seed 42, approximately tied), and 65.52 (seed 123, approximately tied). In two of five scenes under seed 42 and two of five under seed 123, Reflection-only actually exceeds the seed-7 Pure baseline. This pattern suggests that LLM-based self-critique, when applied in isolation, can be helpful, neutral, or harmful depending on the plan initialization trajectory---echoing the single-scenario coastal finding in Section~\ref{sec:results}. The practical implication is that reflection should not be deployed as a standalone improvement mechanism; within the full Memento--Hope--Reflection architecture, memory supplies evidence about what has worked in similar past cases, Hope ensures that capability weights reflect both doctrinal priors and scenario-specific adaptation, and reflection then fine-tunes the plan against the simulator's feedback on a stronger foundation.

The hyperparameter sensitivity analysis (Tables~\ref{tab:hyperparam-theta} and~\ref{tab:hyperparam-epsilon}, recorded under the earlier RTX 3060 environment) reveals a practically important finding about the HOPE mechanism: both SDE parameters ($\theta$ and $\epsilon$) exhibit U-shaped sensitivity curves in that recorded sweep, and the sweep suggested that a systematic hyperparameter optimization pass could unlock additional performance. This conclusion is specific to the recorded sweep; under the 4090 default-parameter refresh, the default Full pipeline (65.76) already exceeds the default Hope-only variant (64.45), so the sweep should be re-run on the current environment before generalizing the finding.

The practical consequence is that the paper should be read not only as a claim of average improvement, but also as a map of where the present pipeline still appears structurally fragile. The multi-seed replication across five seeds demonstrates the Full pipeline's gain is consistent and statistically significant under default HOPE parameters on the RTX 4090 D environment (paired $t(4)=3.75$, $p=0.02$, Cohen's $d=1.68$), with the largest gain in the seed where the baseline was weakest (seed 123, +7.71 points). The 20-iteration convergence run confirms that extended optimization provides further improvement, but breach remains the persistent bottleneck. Just as importantly, the present loop should not be mistaken for a full reinforcement-learning or search regime in the style of standard RL formulations, AlphaGo-scale tree search, or population-based training \citep{sutton2018rl,silver2016alphago,jaderberg2017pbt}. It is closer to a bounded local update rule operating over multiple partially competing mission metrics, a setting that shares some flavor with multi-task optimization trade-offs \citep{yu2020pcgrad}.

\paragraph{Threats to validity.} The five-stage kill-chain simulator is an internal simulation abstraction that has not been calibrated against established military simulation platforms. A parameter sensitivity analysis across 20 perturbation conditions confirmed that the qualitative finding (Full $>$ Pure) remains robust. Readers should interpret reported confidence intervals with care: all CIs in the cross-scenario tables reflect within-seed Monte Carlo variance only (50 runs, seed 7). The cross-seed SD of Pure LLM MS ($\pm1.09$, Table~\ref{tab:multiseed}) is approximately 5$\times$ wider than the within-seed CI widths ($\sim$0.1--0.2 points). The multi-seed replication covers all five seeds under the default Full setting, and the Reflection-only finding in the single-scenario ablation reveals that standalone reflection degrades plan quality under this fixed-seed setting.

Several limitations remain:

\begin{itemize}
\item \textbf{LLM backend traceability}: the refreshed canonical ablation and 20-iteration scene-out evidence now serialize machine-readable runtime manifests. The rerun records identify the local backend model as \path{deepseek-r1-0528-qwen3-8b} and place it in a dedicated Python environment on a 13th Gen Intel(R) Core(TM) i7-13700KF CPU with an NVIDIA GeForce RTX 4090 D GPU (24\,GB), using an OpenAI-compatible local endpoint. This removes the earlier dependence on author-confirmed runtime metadata for the main reported bundles.
\item \textbf{Experimental depth}: the multi-seed replication across five seeds now provides paired $t$-test evidence for the default Full setting ($t(4)=3.75$, $p=0.02$). The three external baselines bracket performance. The Reflection-only single-scenario result shows a clear degradation relative to the pure baseline (59.29 vs.\ 62.75 MS). The per-component ablation (Memento-only, Hope-only) is multi-seed only for seeds 7 and 42. The multi-seed evidence comes from a single scenario family.
	\item \textbf{Reflection-only single-scenario result}: under the fixed seed 7, standalone reflection degrades every headline metric and every stage score relative to the pure baseline (MS 59.29 vs.\ 62.75). This is consistent with the coupled-architecture argument that reflection requires memory-supplied evidence and adaptive weighting to act as a useful repair signal. See Section~\ref{sec:results} and Section~\ref{sec:discussion} for details.
\item \textbf{Case-bank comparability}: the ablation evidence uses a five-record frozen seed bank, whereas the selected generalization evidence uses 250 source cases and 200-case scene-out banks. This is informative, but it also means that the role of memory quality and memory scale must be discussed carefully. In particular, the weak Memento-only result in the ablation table may reflect severe case-bank undersizing (five records) as much as memory-module weakness. A case-bank-size ablation (5/25/100/200 records, Table~\ref{tab:casebank-ablation}, recorded under the earlier RTX 3060 environment) confirms that memory scale matters: the 5-record frozen seed bank is indeed undersized, and the 100-record setting achieves the best overall result.
\item \textbf{Runtime-cost reporting}: the refreshed result bundles now record runtime manifests and wall-clock timing summaries. The cost--quality table provides approximate wall-clock estimates for all settings and baselines. Wall-clock timing data have been integrated into the evaluation platform to reflect resource expenditures.
\item \textbf{Figure scope}: the main-text figures are now generated directly from the system's structured output, and a separate graphical abstract has been prepared as a non-generative submission asset. The remaining work is narrower: final grayscale verification, any journal-specific artwork packaging, and any extra venue-driven figure refinements.
\item \textbf{Metadata completeness}: the author list, affiliation, corresponding-author email, no-funding statement, competing-interest declaration, and CRediT roles are now recorded, but the final acknowledgments wording and data-availability wording may still need journal-specific refinement.

\end{itemize}

These limitations do not invalidate the current contribution. The multi-seed replication for the default Full setting ($n=5$, $p=0.02$, $d=1.68$) provides strong statistical evidence for the Full-pipeline gain, with the overall-effectiveness gain significant at $p<0.01$. The three-backend comparison confirms that the Full pipeline's improvement is not specific to a single LLM. The external baselines bracket performance, and the face validity parameter sensitivity analysis (20/20 perturbation conditions) demonstrates robustness to simulator parameter choices. The next iteration should extend multi-seed evidence to per-component variants (Memento-only, Hope-only) across all generalization scenarios.
